\documentclass[a4paper,12pt]{article}
\usepackage[T2A]{fontenc}
\usepackage[utf8]{inputenc}
\usepackage[russian]{babel}
\usepackage{amsmath,amssymb,mathtools}
\usepackage{helvet}
\renewcommand{\familydefault}{\sfdefault}
\usepackage{icomma}
\usepackage[a4paper,left=20mm,right=20mm,top=20mm,bottom=22mm]{geometry}
\usepackage{microtype,setspace,parskip,xcolor,enumitem,booktabs,tabularx,fancyhdr,hyperref}
\usepackage{tikz}
\usetikzlibrary{calc,arrows.meta,positioning,shapes.geometric}
\usepackage{pgfplots}
\pgfplotsset{compat=1.18}
\onehalfspacing
\definecolor{accent}{HTML}{008080}
\definecolor{darkaccent}{HTML}{004D4D}
\definecolor{textgray}{HTML}{333333}
\definecolor{softgray}{HTML}{6E7777}
\color{textgray}
\pagestyle{fancy}\fancyhf{}\fancyhead[L]{\small\textcolor{softgray}{ОГЭ · математика}}\fancyhead[R]{\small\textcolor{softgray}{Марина Селиверстова}}\fancyfoot[C]{\small\textcolor{softgray}{\thepage}}\renewcommand{\headrulewidth}{0pt}
\setlist[itemize]{leftmargin=6mm,itemsep=2pt,topsep=4pt}
\setlist[enumerate]{leftmargin=8mm,itemsep=4pt,topsep=4pt}
\begin{document}
\begin{titlepage}\thispagestyle{empty}\centering
\vspace*{24mm}{\large\textcolor{softgray}{Математика · ОГЭ}\par}\vspace{17mm}
{\fontsize{27}{32}\selectfont\bfseries\color{darkaccent}Чек-лист\par}\vspace{7mm}
{\Large Парабола, прямая $y=m$\\[2mm] и графики дробно-рациональных функций\par}\vspace{8mm}
{\large ОДЗ · разложение на множители · выколотая точка\par}\vspace{22mm}
{\large Марина Селиверстова\par}\vspace{7mm}{\normalsize 16.09.2026\par}\vfill
{\normalsize Лёвин Артём Александрович\\эксклюзивно для Марины Селиверстовой\par}\vspace{17mm}
\begin{tikzpicture}\draw[accent,line width=1pt](-6.2,0)..controls(-3.8,.7)and(-1.7,-.65)..(0,0)..controls(2,.75)and(4,-.55)..(6.2,.1);\end{tikzpicture}
\end{titlepage}

\section*{Главная идея}
В задачах с графиками удобно разделять работу на две части: привести функцию к виду, который легко строить, и сохранить ограничения исходного выражения. Для параболы
\[y=ax^2+bx+c,\qquad a\ne0\]
основные ориентиры: вершина, пересечения с осями $Ox$ и $Oy$, направление ветвей. Если исходная функция содержит дробь, после сокращения сохраняется ОДЗ исходной функции. Исключённое значение аргумента изображается выколотой точкой.

\section*{Парабола: что нужно помнить}
\[
 a>0\Rightarrow\text{ветви вверх},\qquad a<0\Rightarrow\text{ветви вниз}.
\]
Абсцисса вершины:
\[\boxed{x_{\mathrm{в}}=-\frac{b}{2a}}.\]
Ордината вершины:
\[\boxed{y_{\mathrm{в}}=f(x_{\mathrm{в}})}.\]
Полезное обратное соотношение:
\[\boxed{b=-2ax_{\mathrm{в}}}.\]
Пересечение с $Oy$ находится при $x=0$, поэтому для $y=ax^2+bx+c$ это точка $(0;c)$. Пересечения с $Ox$ находятся из уравнения $f(x)=0$.

\section*{Пример 1. Построение $y=x^2-x-6$}
Здесь $a=1$, $b=-1$, $c=-6$.
\[
x_{\mathrm{в}}=-\frac{-1}{2}=\frac12,
\qquad
y_{\mathrm{в}}=\left(\frac12\right)^2-\frac12-6=-\frac{25}{4}=-6{,}25.
\]
Вершина:
\[\boxed{V\left(\frac12;-\frac{25}{4}\right)}.\]
Пересечения с $Ox$:
\[x^2-x-6=0\iff(x-3)(x+2)=0,\]
поэтому
\[\boxed{A(-2;0),\qquad B(3;0)}.\]
Пересечение с $Oy$:
\[\boxed{C(0;-6)}.\]
Для экзаменационного эскиза достаточно корректно отметить вычисленные точки и плавно провести параболу.

\section*{Прямая $y=m$}
Уравнение $y=m$ задаёт горизонтальную прямую. Количество решений $f(x)=m$ совпадает с количеством общих точек графиков $y=f(x)$ и $y=m$.

Для $y=x^2-x-6$ минимальное значение равно $-25/4$, поэтому
\[
\boxed{m<-\frac{25}{4}\Rightarrow0\text{ точек}},
\]
\[
\boxed{m=-\frac{25}{4}\Rightarrow1\text{ точка}},
\]
\[
\boxed{m>-\frac{25}{4}\Rightarrow2\text{ точки}}.
\]

\section*{Сложная задача: сначала упростить}
Рассмотрим
\[
\boxed{y=\frac{(x+4)(x^2+3x+2)}{x+1}}.
\]
Сначала выписываем ОДЗ:
\[x+1\ne0\Rightarrow\boxed{x\ne-1}.\]
Квадратный трёхчлен раскладываем на множители:
\[x^2+3x+2=(x+1)(x+2).\]
Тогда при $x\ne-1$
\[
y=\frac{(x+4)(x+1)(x+2)}{x+1}=(x+4)(x+2)=x^2+6x+8.
\]
Ограничение $x\ne-1$ сохраняется.

Для упрощённой параболы:
\[
x_{\mathrm{в}}=-3,\qquad y_{\mathrm{в}}=-1,
\]
то есть $V(-3;-1)$. Нули функции: $x=-4$ и $x=-2$, пересечение с $Oy$: $(0;8)$.

Чтобы найти выколотую точку, подставляем исключённое значение $x=-1$ в упрощённую формулу:
\[
y=(-1)^2+6(-1)+8=3.
\]
Следовательно,
\[\boxed{Q(-1;3)}\]
--- выколотая точка.

\section*{Ещё один пример}
\[
y=\frac{(x+1)(x^2+7x+10)}{x+2}.
\]
ОДЗ: $x\ne-2$. Так как $x^2+7x+10=(x+2)(x+5)$, то
\[
y=(x+1)(x+5)=x^2+6x+5,\qquad x\ne-2.
\]
При $x=-2$ получаем $y=-3$, поэтому $Q(-2;-3)$ --- выколотая точка.

\section*{Алгоритм}
\begin{enumerate}
\item Найдите ОДЗ исходной функции.
\item Разложите квадратные трёхчлены на множители, если это возможно.
\item Сократите общий множитель, сохранив исходное ОДЗ.
\item Получите удобную формулу и найдите вершину, пересечения с $Ox$ и $Oy$.
\item Для исключённого значения $x$ вычислите соответствующее $y$ по упрощённой формуле.
\item Постройте график и отметьте выколотую точку пустым кружком.
\end{enumerate}

\section*{Типичные ошибки}
\begin{enumerate}
\item Потерять ОДЗ после сокращения.
\item Перепутать вершину и пересечение с осью $Oy$.
\item При пересечении с $Ox$ обнулить $x$ вместо $y$.
\item Ошибиться в арифметике при вычислении вершины.
\item Неправильно извлечь $\sqrt D$.
\item Сокращать слагаемые: сокращение возможно только для множителей.
\item Не подписать ключевые и выколотые точки.
\end{enumerate}

\clearpage
\section*{Тренировочный блок}
\begin{enumerate}
\item Для $y=x^2+2x-3$ найдите вершину, пересечения с $Ox$ и $Oy$. Постройте эскиз.
\item Для $y=x^2-4x+3$ определите, при каких $m$ прямая $y=m$ имеет с параболой 0, 1 и 2 общие точки.
\item Постройте график $\displaystyle y=\frac{(x+3)(x^2+x-2)}{x+2}$. Укажите ОДЗ и выколотую точку.
\item Постройте график $\displaystyle y=\frac{(x-2)(x^2-5x+6)}{x-3}$ и найдите все $m$, при которых $y=m$ имеет с графиком ровно одну общую точку.
\item Для $\displaystyle y=\frac{(x+1)(x^2-4x+3)}{x-1}$ найдите ОДЗ, упростите выражение, постройте график с выколотой точкой и определите число общих точек с $y=m$ в зависимости от $m$.
\end{enumerate}

\clearpage
\section*{Ответы}
\renewcommand{\arraystretch}{1.35}
\begin{tabularx}{\linewidth}{p{10mm}X}\toprule
\textbf{№}&\textbf{Ответ}\\\midrule
1&$V(-1;-4)$; $Ox$: $(-3;0)$, $(1;0)$; $Oy$: $(0;-3)$.\\
2&$m<-1$: 0 точек; $m=-1$: 1 точка; $m>-1$: 2 точки.\\
3&$x\ne-2$; $y=x^2+2x-3$; $V(-1;-4)$; нули $-3$ и $1$; $Q(-2;-3)$.\\
4&$x\ne3$; $y=(x-2)^2$; $Q(3;1)$; ровно одна общая точка при $m=0$ или $m=1$.\\
5&$x\ne1$; $y=(x-1)^2-4$; $Q(1;-4)$; при $m\le-4$ общих точек нет, при $m>-4$ две общие точки.\\\bottomrule
\end{tabularx}
\end{document}
